\scanpage{56}
\mspage{55}

c) $\varnothing\in\Phi$ \deleted{$\ne\varnothing$ [en vertu de b), c'est tjrs vrai si $ M\ne\varnothing$], et si $ M=\varnothing$, cela signifie que $\varnothing$}

\deleted{\uline{\unsure{Rép.|Rq.}} (équivaut : $\varnothing\in\Phi$)}

\uline{Remarque} Inversement, soient $ M$ un ens, $\deleted{\Phi}\mathcal F\subset\mathcal P( M)$ satisfaisant les conditions a) b) c). NB La condition c résulte de \inserted{a),} b si $ M\ne\varnothing$. Dans tous les cas, elle équivaut à : \deleted{\unsure{$\in$|à}} c') $\deleted{\Phi}\mathcal F\ne\varnothing$ i.e existe $\deleted{X}A\subset M$ telle que $A\in\deleted{\Phi}\mathcal F$

% slanted
\margin{\unsure{Ces|Les} conditions impliquent que $\mathcal F$ est stable par $A\cap B$ (mais en général \unsure{on n'a|\ldots} $\not\Rightarrow$ $ M\in\mathcal F$ !)}

Alors

1°) $\Phi$ ordonné par inclusion \deleted{\unsure{ad|sa}}satisfait les conditions C1, C2. \deleted{Les objets irréductibles} \deleted{\unsure{sont}}

2°) Les objets irréductibles sont les $A_X$, et $X\mapsto\deleted{X}A_X$ est une bijection de $ M$ \deleted{avec} l'ens des objets irred. de $\deleted{\Phi}\mathcal F$.

3°) Si $X\in M$, $A\in\deleted{\Phi}\mathcal F$, alors $X\in A\Longleftrightarrow A_X\subset A$ \deleted{, donc} (tautologique) donc \deleted{$A\in\mathcal F$} si on identifie $ M$ \inserted{(par $X\unsure{\vdash|\mapsto} A_X$)} à l'ens des objets irréductibles de $\deleted{\Phi}\mathcal F$, alors pour $A\in\mathcal F$, $\overline A=\{A_X\mid A_X\subset A\inserted{\text{ie.}\,X\in A}\}$ s'identifie à $A$ lui-$=$.

\centerline{\rule{1.8cm}{0.35pt}}

On considérera $ M$ comme ordonné par $\leq$ (jamais par $\prec$ !) : multistrates et \uline{sous-multistrates}. Les figures \inserted{$F$} peuvent s'identifier aux \deleted{des parties} ens. de multistrates $\overline F$ correspondants (\unsure{ch. que|cl. par|chaque} figure à l'ens de ses strates).
